Il lavoro ha formulato e implementato un modello ridotto Biot--Darcy con una legge di Kelvin--Voigt frazionaria per la matrice e un tracciante passivo con accumulo frazionario. Le verifiche MMS mostrano errori decrescenti e ordini spaziali osservati prossimi a quelli attesi per lo spostamento in $L^2$ e per pressione e concentrazione; tuttavia, la serie coordinata non isola completamente l'errore spaziale del tracciante e non sostituisce un'analisi teorica completa dell'errore né una verifica inf--sup. I test di massa raggiungono drift dell'ordine della precisione macchina, mentre l'identità algebrica passo-passo e il residuo energetico cumulativo sono stati trattati come controlli distinti.

La mesh quadrilatera derivata da una sezione BraTS 2020 riproduce con elevata sovrapposizione la maschera adottata, ma il dominio è un proxy del parenchima, non una segmentazione tissutale specifica della sostanza bianca. Le sorgenti e i campi temporali della dimostrazione anatomica sono sintetici; non sono disponibili osservazioni cliniche longitudinali che consentano di validare le traiettorie o di interpretare le stime come parametri fisiologici.

L'analisi del Jacobiano discreto conferma la degenerazione prevista quando il flusso è nullo e mostra rango pieno locale nel protocollo attivo. I profili ottimizzati indicano che il valore preferito di $\beta$ dipende dalla componente del funzionale considerata: il minimo campionato del termine dati è in $0.45$, mentre quello del funzionale regolarizzato è in circa $0.479$. Il profilo dati è inoltre sensibile alla discretizzazione inversa: passando da $(N,M)=(16,80)$ a $(24,120)$ il minimo campionato si sposta a $0.4375$ e il valore minimo del profilo diminuisce. Queste osservazioni sono coerenti con una valle di costo piatta e con un effetto non trascurabile del prior, ma non dimostrano da sole non-identificabilità strutturale di $\beta$ nel protocollo attivo.

La campagna finale comprende tre run senza rumore e dieci run rumorose con budget pianificato di 600 valutazioni FEM. La dispersione delle stime è contenuta per alcuni parametri nei campioni osservati, ma resta marcata per $\beta$ ed $E$; gli errori non variano monotonicamente con il livello nominale di rumore. Tutte le esecuzioni hanno esaurito il budget senza certificare convergenza, e le dieci run rumorose condividono lo stesso seed CPSO. I dati supportano quindi una caratterizzazione controllata del problema inverso sintetico, non una dimostrazione generale di robustezza dell'algoritmo.

Prima della sottomissione occorre completare almeno quattro verifiche: (i) estendere la sensibilità della discretizzazione e del passo temporale, separando gli errori del forward solver da quelli di ottimizzazione; (ii) aumentare le repliche e incrociare seed del rumore e seed CPSO per quantificare entrambe le fonti di variabilità; (iii) confrontare eventuali ottimizzatori alternativi a parità di funzione obiettivo, discretizzazione e budget, se si intendono formulare affermazioni comparative; (iv) depositare un archivio versionato che permetta di rigenerare ogni tabella e figura. Un'applicazione clinica alla sostanza bianca richiederebbe inoltre una segmentazione tissutale indipendente e osservazioni adeguate per la validazione. In assenza di tali dati, la conclusione appropriata è quella di una verifica numerica e metodologica preliminare, con limiti esplicitati.
